\chapter{State, roadmap and open decisions}
\label{ch:roadmap}

\section{What exists}

On branch \texttt{amr-complex}, all suites passing in two and three dimensions
at one to eight ranks:

\begin{table}[htbp]
\centering
\caption{Components added.}
\label{tab:files}
\begin{tabular}{p{0.34\textwidth}p{0.58\textwidth}}
\toprule
file & role \\
\midrule
\texttt{Src/AMRComplex.\{H,cpp\}} &
  conforming multilevel complex; identity $(\lev,\iv)$ \\
\texttt{Src/AMRContour.\{H,cpp\}} &
  marching simplices, oriented, carried components, MEF \\
\texttt{Src/AMRSegment.\{H,cpp\}} &
  connected components of a threshold set \\
\texttt{Src/AMRParallel.\{H,cpp\}} &
  the only place MPI calls appear \\
\texttt{Src/amrComplexVerify.cpp} &
  complex gate \\
\texttt{Src/amrContourVerify.cpp} &
  surface gate, negative controls \\
\texttt{Src/amrSweep.cpp} &
  threshold sweep \\
\texttt{Src/makeTestPlt.cpp} &
  synthetic hierarchies with analytic references \\
\texttt{Src/GNUmakefile.AMRComplex} &
  standalone, AMReX only, \texttt{DIM ?= 3} \\
\texttt{Tools/amr*.sh} &
  the three harnesses \\
\bottomrule
\end{tabular}
\end{table}

\section{Roadmap}

\begin{enumerate}[leftmargin=*]
\item \textbf{Representative snapshots.} One plotfile carrying two or three
  candidate surrogates---heat release rate, temperature, an
  intermediate---swept individually, with the $N(v)$ curves laid side by side.
  This answers whether a plateau exists, supplies real timings for
  \texttt{ownedBaseRegions} and the segmentation maps at production size, and
  shows whether the surrogates agree, all before anything is built on an
  assumption about any of the three.
\item \textbf{Segmentation as a plotfile}, and contour component counts.
\item \textbf{Merge tree}, gated on \cref{rule:sweepgate}, and subject to the
  decision of \cref{ch:gpu}.
\item \textbf{Overlap matching} (\cref{sec:overlap}), serving both time tracking
  and cross-surrogate comparison.
\item \textbf{Boundary invariants} for the $\betti{1}$ events.
\end{enumerate}

\section{Open decisions}

\decide{Whether a cell is a region thresholded from above or from below. It may
differ per surrogate, so it remains a runtime choice rather than a design
commitment.}

\decide{Host or device for the merge tree (\cref{ch:gpu}).}

\decide{Wall extrapolation. Cell-centred data leaves the complex inset by
$\Delta x/2$ at each non-periodic wall. Zero normal gradient is a defensible
placeholder and is exactly right for a symmetry, slip or adiabatic wall, but it
is exactly wrong at a Dirichlet wall, and it introduces a prismatic collar whose
area contribution is $P\,\Delta x/2$ for inset-plane perimeter $P$, together
with a spurious curvature ridge standing off each wall at a fixed
$\Delta x/2$. The error is bounded by how obliquely the surface meets the wall,
so the incidence-angle distribution on the inset plane decides whether it
matters at all. Limited linear extrapolation is the recommended default; a
supplied boundary-condition table is the only correct option at a Dirichlet
wall.}

\section{Known limits}

\begin{itemize}[leftmargin=*]
\item MEF output gathers to the I/O processor---the same wall
  \texttt{isosurface.cpp} has. \texttt{.pvtu} is the intended replacement.
\item \texttt{ownedBaseRegions} and the segmentation maps have not been exercised
  at production size.
\item All verification to date uses synthetic fields on generated hierarchies.
  Hierarchies from actual regridding will exercise the proper-nesting assertion
  of \cref{sec:stitch} in a way synthetic ones cannot; if it aborts, that is the
  check working.
\item Host-only, per \cref{ch:gpu}.
\end{itemize}
